% problem_id: S001-artin-lemma-approximate
% source_id: S001 (Stacks Project)
% source_file: artin.tex
% statement_locator: artin.tex lines 1270--1292
% proof_locator: artin.tex lines 1294--1432
% env: lemma
% id_note: label 跨文件重名 ⇒ id 用文件限定形式
% pairing: tight (verified)
% status: complete (statement + author proof verbatim + reason + locators)
%
\begin{lemma}
\label{lemma-approximate}
Let $S$ be a locally Noetherian scheme. Let
$p : \mathcal{X} \to (\Sch/S)_{fppf}$ be a category
fibred in groupoids. Let $x$ be an object of
$\mathcal{X}$ lying over $\Spec(R)$ where $R$ is a Noetherian complete
local ring with residue field $k$ of finite type over $S$. Let $s \in S$
be the image of $\Spec(k) \to S$. Assume that (a) $\mathcal{O}_{S, s}$ is
a G-ring and (b) $p$ is limit preserving on objects. Then for every
integer $N \geq 1$ there exist
\begin{enumerate}
\item a finite type $S$-algebra $A$,
\item a maximal ideal $\mathfrak m_A \subset A$,
\item an object $x_A$ of $\mathcal{X}$ over $\Spec(A)$,
\item an $S$-isomorphism $R/\mathfrak m_R^N \cong A/\mathfrak m_A^N$,
\item an isomorphism
$x|_{\Spec(R/\mathfrak m_R^N)} \cong x_A|_{\Spec(A/\mathfrak m_A^N)}$
compatible with (4), and
\item an isomorphism
$\text{Gr}_{\mathfrak m_R}(R) \cong \text{Gr}_{\mathfrak m_A}(A)$
of graded $k$-algebras.
\end{enumerate}
\end{lemma}

\begin{proof}
Choose an affine open $\Spec(\Lambda) \subset S$ such that $k$ is a finite
$\Lambda$-algebra, see
Morphisms, Lemma \ref{morphisms-lemma-point-finite-type}.
We may and do replace $S$ by $\Spec(\Lambda)$.

\medskip\noindent
We may write $R$ as a directed colimit $R = \colim C_j$ where each
$C_j$ is a finite type $\Lambda$-algebra (see
Algebra, Lemma \ref{algebra-lemma-ring-colimit-fp}).
By assumption (b) the object $x$ is isomorphic to the restriction of
an object over one of the $C_j$. Hence we may choose a finite type
$\Lambda$-algebra $C$, a $\Lambda$-algebra map $C \to R$, and an object
$x_C$ of $\mathcal{X}$ over $\Spec(C)$ such that $x = x_C|_{\Spec(R)}$.
The choice of $C$ is a bookkeeping device and could be avoided.
For later use, let us write $C = \Lambda[y_1, \ldots, y_u]/(f_1, \ldots, f_v)$
and we denote $\overline{a}_i \in R$ the image of $y_i$ under the
map $C \to R$. Set $\mathfrak m_C = C \cap \mathfrak m_R$.

\medskip\noindent
Choose a $\Lambda$-algebra surjection $\Lambda[x_1, \ldots, x_s] \to k$
and denote $\mathfrak m'$ the kernel.
By the universal property of polynomial rings we may lift this
to a $\Lambda$-algebra map $\Lambda[x_1, \ldots, x_s] \to R$.
We add some variables (i.e., we increase $s$ a bit) mapping to generators
of $\mathfrak m_R$. Having done this we see that
$\Lambda[x_1, \ldots, x_s] \to R/\mathfrak m_R^2$ is surjective.
Then we see that
\begin{equation}
\label{equation-surjection}
P = \Lambda[x_1, \ldots, x_s]_{\mathfrak m'}^\wedge \longrightarrow R
\end{equation}
is a surjective map of Noetherian complete local rings, see for example
Formal Deformation Theory, Lemma
\ref{formal-defos-lemma-surjective-cotangent-space}.

\medskip\noindent
Choose lifts $a_i \in P$ of $\overline{a}_i$ we found above.
Choose generators $b_1, \ldots, b_r \in P$ for the kernel of
(\ref{equation-surjection}).
Choose $c_{ji} \in P$ such that
$$
f_j(a_1, \ldots, a_u) = \sum c_{ji} b_i
$$
in $P$ which is possible by the choices made so far. Choose generators
$$
k_1, \ldots, k_t \in
\Ker(P^{\oplus r} \xrightarrow{(b_1, \ldots, b_r)} P)
$$
and write $k_i = (k_{i1}, \ldots, k_{ir})$ and $K = (k_{ij})$
so that
$$
P^{\oplus t} \xrightarrow{K}
P^{\oplus r} \xrightarrow{(b_1, \ldots, b_r)}
P \to R \to 0
$$
is an exact sequence of $P$-modules. In particular we have
$\sum k_{ij} b_j = 0$. After possibly increasing $N$ we may
assume $N - 1$ works in the Artin-Rees lemma for the first two maps of this
exact sequence (see More on Algebra, Section
\ref{more-algebra-section-artin-rees} for terminology).

\medskip\noindent
By assumption $\mathcal{O}_{S, s} = \Lambda_{\Lambda \cap \mathfrak m'}$ is
a G-ring. Hence by More on Algebra, Proposition
\ref{more-algebra-proposition-finite-type-over-G-ring}
the ring $\Lambda[x_1, \ldots, x_s]_{\mathfrak m'}$ is a $G$-ring.
Hence by Smoothing Ring Maps, Theorem
\ref{smoothing-theorem-approximation-property-variant}
there exist an \'etale ring map
$$
\Lambda[x_1, \ldots, x_s]_{\mathfrak m'} \to B,
$$
a maximal ideal $\mathfrak m_B$ of $B$ lying over $\mathfrak m'$, and
elements $a'_i, b'_i, c'_{ij}, k'_{ij} \in B'$ such that
\begin{enumerate}
\item $\kappa(\mathfrak m') = \kappa(\mathfrak m_B)$ which implies
that $\Lambda[x_1, \ldots, x_s]_{\mathfrak m'} \subset B_{\mathfrak m_B}
\subset P$ and $P$ is identified with the completion of $B$ at
$\mathfrak m_B$, see remark preceding Smoothing Ring Maps, Theorem
\ref{smoothing-theorem-approximation-property-variant},
\item $a_i - a'_i, b_i - b'_i, c_{ij} - c'_{ij}, k_{ij} - k'_{ij} \in
(\mathfrak m')^N P$, and
\item $f_j(a'_1, \ldots, a'_u) = \sum c'_{ji} b'_i$ and $\sum k'_{ij}b'_j = 0$.
\end{enumerate}
Set $A = B/(b'_1, \ldots, b'_r)$ and denote $\mathfrak m_A$ the
image of $\mathfrak m_B$ in $A$. (Note that $A$ is essentially of finite
type over $\Lambda$; at the end of the proof we will show how to obtain
an $A$ which is of finite type over $\Lambda$.) There is a ring map
$C \to A$ sending $y_i \mapsto a'_i$ because the $a'_i$ satisfy
the desired equations modulo $(b'_1, \ldots, b'_r)$.
Note that $A/\mathfrak m_A^N = R/\mathfrak m_R^N$ as quotients of
$P = B^\wedge$ by property (2) above. Set $x_A = x_C|_{\Spec(A)}$.
Since the maps
$$
C \to A \to A/\mathfrak m_A^N \cong R/\mathfrak m_R^N
\quad\text{and}\quad
C \to R \to R/\mathfrak m_R^N
$$
are equal we see that $x_A$ and $x$ agree modulo $\mathfrak m_R^N$
via the isomorphism $A/\mathfrak m_A^N = R/\mathfrak m_R^N$. At this
point we have shown properties (1) -- (5) of the statement of the lemma.
To see (6) note that
$$
P^{\oplus t} \xrightarrow{K}
P^{\oplus r} \xrightarrow{(b_1, \ldots, b_r)}
P
\quad\text{and}\quad
P^{\oplus t} \xrightarrow{K'}
P^{\oplus r} \xrightarrow{(b'_1, \ldots, b'_r)}
P
$$
are two complexes of $P$-modules which are congruent modulo
$(\mathfrak m')^N$ with the first one being exact. By our choice of $N$
above we see from
More on Algebra, Lemma \ref{more-algebra-lemma-approximate-complex-graded}
that $R = P/(b_1, \ldots, b_r)$ and
$P/(b'_1, \ldots, b'_r) = B^\wedge/(b'_1, \ldots, b'_r) = A^\wedge$
have isomorphic associated graded algebras, which is what we wanted to show.

\medskip\noindent
This last paragraph of the proof serves to clean up the issue that $A$ is
essentially of finite type over $S$ and not yet of finite type.
The construction above gives $A = B/(b'_1, \ldots, b'_r)$ and
$\mathfrak m_A \subset A$ with $B$ \'etale over
$\Lambda[x_1, \ldots, x_s]_{\mathfrak m'}$. Hence $A$ is of finite
type over the Noetherian ring $\Lambda[x_1, \ldots, x_s]_{\mathfrak m'}$.
Thus we can write $A = (A_0)_{\mathfrak m'}$ for some finite type
$\Lambda[x_1, \ldots, x_n]$ algebra $A_0$. Then
$A = \colim (A_0)_f$ where
$f \in \Lambda[x_1, \ldots, x_n] \setminus \mathfrak m'$, see
Algebra, Lemma \ref{algebra-lemma-localization-colimit}.
Because $p : \mathcal{X} \to (\Sch/S)_{fppf}$ is limit preserving on
objects, we see that
$x_A$ comes from some object $x_{(A_0)_f}$ over $\Spec((A_0)_f)$ for
an $f$ as above. After replacing $A$ by $(A_0)_f$ and $x_A$ by
$x_{(A_0)_f}$ and $\mathfrak m_A$ by $(A_0)_f \cap \mathfrak m_A$
the proof is finished.
\end{proof}
